%==============================================================================
%  04-theoretical-foundations.tex — Theoretical Foundations
%==============================================================================
\strategyPart{04 \textbar\ Theory}{Theoretical Foundations}

Five established theories anchor the survey. Together they explain why AI
transformation is an organizational problem, not a model-performance problem.

\section{The five lenses}

\vspace{-0.3em}
\infobox{\textbf{1 · Digital transformation theory}}{%
  Organizations use digital technologies to change value creation, processes,
  customer relationships and operating models [MIT Sloan; HBR; Gartner;
  Deloitte]. AI is the newest wave of this phenomenon.}

\infobox{\textbf{2 · Sociotechnical systems theory}}{%
  AI is a joint system of technology, people, work practices, structures and
  rules. Performance depends on how these co-evolve — not on model quality
  alone [academic STS literature; HCI research].}

\infobox{\textbf{3 · Resource-based view (RBV)}}{%
  Advantage comes from valuable, rare, inimitable resources: proprietary data,
  AI talent, domain expertise, infrastructure, governance and trust
  [strategic-management literature; McKinsey; Deloitte].}

\infobox{\textbf{4 · Dynamic capabilities}}{%
  Advantage is a capacity to \textit{sense} AI opportunities, \textit{seize}
  them through investment, and \textit{transform} structures continuously
  [Teece et al.; dynamic-capabilities literature].}

\infobox{\textbf{5 · Human--AI collaboration theory}}{%
  Humans and AI divide responsibility along a spectrum: AI informs, AI proposes
  and human approves, AI executes and human monitors, or agents execute under
  human governance [HBR 2018; NIST AI RMF; Microsoft; WEF].}

\section{What the theories imply}

\begin{itemize}
  \item RBV and dynamic capabilities explain \textbf{why the gap between
        adoption and scaling persists}: tools are commodity resources; the
        stocks around them are not [MIT SMR--BCG 2019].
  \item Sociotechnical theory explains \textbf{why AI failures are
        organizational}: a technically excellent model embedded in a
        poorly designed workflow still fails.
  \item Human--AI collaboration theory explains \textbf{why adoption is the
        bottleneck}: 30\% of generative-AI projects are abandoned at the
        point where humans must begin working with the system [Gartner 2023].
\end{itemize}
